% GENERATED FILE. Edit canonical lessons, then run npm run generate:books.
% canonical-source-hash: fnv1a64:ca6bab859f298720
% canonical-lessons: RU-C401-prilozhenie, RU-C401-ustanavlivat, RU-C401-obnovlyat, RU-C401-udalyat, RU-C401-zagruzhat

\chapter{Apps}
\label{ch:ru-apps}
\hlchaptermodality{401}

\begin{chapteropening}
\textbf{By the end of this chapter:} \emph{I can talk about apps: installing them, updating them, deleting them and uploading files.}

Complete the last lesson of chapter 401: I can talk about apps: installing them, updating them, deleting them and uploading files.
\end{chapteropening}

\section[\texorpdfstring{\ru{приложение} (prilozhéniye)}{приложение (prilozhéniye)}]{\ru{приложение} (prilozhéniye) — an app; an attachment; a supplement}

\label{lesson:RU-C401-prilozhenie}



\begin{quote}
Before the new one: say the Russian for to break into, then the Russian for a virus.
\end{quote}

\subsection*{You'll want to know: \ru{приложение}}
\textbf{\ru{приложение}} (\emph{prilozhéniye}) — \textquotedblleft{}an app; an attachment; a supplement\textquotedblright{}.

\begin{quote}
\ru{Я} \ru{скачал} \ru{новое} \ru{приложение}.
\end{quote}

\begin{quote}
\emph{Ya skachál nóvoye prilozhéniye.}
\end{quote}

\begin{quote}
I downloaded a new app.
\end{quote}

Neuter. An app on a phone, an attachment to a letter, or a supplement to a magazine.

\begin{cousinweb}
\textbf{\ru{приложение}} is built on \emph{\ru{приложить}} \textquotedblleft{}to put to, to attach\textquotedblright{}, from \emph{\ru{при}-} \textquotedblleft{}to\textquotedblright{} and \emph{\ru{ложить}} \textquotedblleft{}to lay\textquotedblright{}: something laid alongside. An app is the newest sense.
\end{cousinweb}

\begin{grammarlens}[title={What you've built}]
One more word for telling, explaining and describing at B1.
\end{grammarlens}

\subsection*{Your turn}
\begin{itemize}
\raggedright
  \item \emph{Say it:} \emph{prilozhéniye}
  \item \emph{Say it:} \emph{prilozhéniye}, once more
  \item \emph{Say it:} a sentence of your own with \emph{prilozhéniye}
  \item \emph{Recall:} say the Russian for to break into, then the Russian for a virus, then say \emph{prilozhéniye} again
\end{itemize}

\begin{tcolorbox}[breakable,colback=teal!4,colframe=teal!35!black,title={Before you move on}]
What does \ru{приложение} mean? (\textquotedblleft{}An app; an attachment; a supplement\textquotedblright{}.) Say it once more.
\end{tcolorbox}
\section[\texorpdfstring{\ru{устанавливать} (ustanávlivat)}{устанавливать (ustanávlivat)}]{\ru{устанавливать} (ustanávlivat) — to install; to establish}

\label{lesson:RU-C401-ustanavlivat}



\begin{quote}
Before the new one: say the Russian for a virus, then the Russian for an app.
\end{quote}

\subsection*{You'll want to know: \ru{устанавливать}}
\textbf{\ru{устанавливать}} (\emph{ustanávlivat}) — \textquotedblleft{}to install; to establish\textquotedblright{}.

\begin{quote}
\ru{Он} \ru{устанавливает} \ru{программу} \ru{на} \ru{компьютер}.
\end{quote}

\begin{quote}
\emph{On ustanávlivayet prográmmu na kompyúter.}
\end{quote}

\begin{quote}
He is installing a program on the computer.
\end{quote}

Imperfective; perfective \emph{\ru{установить}} (\emph{\ru{установлю}, \ru{установишь}}). It also means to set up or establish: \emph{\ru{установить} \ru{рекорд}}.

\begin{cousinweb}
\textbf{\ru{устанавливать}} is built from \emph{\ru{у}-} and \emph{\ru{ставить}} \textquotedblleft{}to put, to stand up\textquotedblright{}, on the root of \emph{\ru{стоять}}, Proto-Slavic \emph{*stojati} (EtymDB): to set a thing firmly in place.
\end{cousinweb}

\begin{grammarlens}[title={What you've built}]
One more word for telling, explaining and describing at B1.
\end{grammarlens}

\subsection*{Your turn}
\begin{itemize}
\raggedright
  \item \emph{Say it:} \emph{ustanávlivat}
  \item \emph{Say it:} \emph{ustanávlivat}, once more
  \item \emph{Say it:} a sentence of your own with \emph{ustanávlivat}
  \item \emph{Recall:} say the Russian for a virus, then the Russian for an app, then say \emph{ustanávlivat} again
\end{itemize}

\begin{tcolorbox}[breakable,colback=teal!4,colframe=teal!35!black,title={Before you move on}]
What does \ru{устанавливать} mean? (\textquotedblleft{}To install; to establish\textquotedblright{}.) Say it once more.
\end{tcolorbox}
\section[\texorpdfstring{\ru{обновлять} (obnovlyát)}{обновлять (obnovlyát)}]{\ru{обновлять} (obnovlyát) — to update, to renew}

\label{lesson:RU-C401-obnovlyat}



\begin{quote}
Before the new one: say the Russian for an app, then the Russian for to install.
\end{quote}

\subsection*{You'll want to know: \ru{обновлять}}
\textbf{\ru{обновлять}} (\emph{obnovlyát}) — \textquotedblleft{}to update, to renew\textquotedblright{}.

\begin{quote}
\ru{Телефон} \ru{сам} \ru{обновляет} \ru{приложения}.
\end{quote}

\begin{quote}
\emph{Telefón sam obnovlyáyet prilozhéniya.}
\end{quote}

\begin{quote}
The phone updates its apps by itself.
\end{quote}

Imperfective; the perfective is \emph{\ru{обновить}} (\emph{\ru{обновлю}, \ru{обновишь}}). \emph{\ru{Обновление}} — an update.

\begin{cousinweb}
\textbf{\ru{обновлять}} is built from \emph{\ru{об}-} and \emph{\ru{новый}} \textquotedblleft{}new\textquotedblright{}, inherited from Proto-Slavic \emph{*novŭ}: to make a thing new again.
\end{cousinweb}

\begin{grammarlens}[title={What you've built}]
One more word for telling, explaining and describing at B1.
\end{grammarlens}

\subsection*{Your turn}
\begin{itemize}
\raggedright
  \item \emph{Say it:} \emph{obnovlyát}
  \item \emph{Say it:} \emph{obnovlyát}, once more
  \item \emph{Say it:} a sentence of your own with \emph{obnovlyát}
  \item \emph{Recall:} say the Russian for an app, then the Russian for to install, then say \emph{obnovlyát} again
\end{itemize}

\begin{tcolorbox}[breakable,colback=teal!4,colframe=teal!35!black,title={Before you move on}]
What does \ru{обновлять} mean? (\textquotedblleft{}To update, to renew\textquotedblright{}.) Say it once more.
\end{tcolorbox}
\section[\texorpdfstring{\ru{удалять} (udalyát)}{удалять (udalyát)}]{\ru{удалять} (udalyát) — to delete, to remove}

\label{lesson:RU-C401-udalyat}



\begin{quote}
Before the new one: say the Russian for to install, then the Russian for to update.
\end{quote}

\subsection*{You'll want to know: \ru{удалять}}
\textbf{\ru{удалять}} (\emph{udalyát}) — \textquotedblleft{}to delete, to remove\textquotedblright{}.

\begin{quote}
\ru{Я} \ru{удаляю} \ru{старые} \ru{фотографии}.
\end{quote}

\begin{quote}
\emph{Ya udalyáyu stáryye fotográfii.}
\end{quote}

\begin{quote}
I am deleting old photos.
\end{quote}

Imperfective; perfective \emph{\ru{удалить}}. It also means to take out a tooth, or to send a player off.

\begin{cousinweb}
\textbf{\ru{удалять}} is built from \emph{\ru{у}-} and the root of \emph{\ru{даль}} \textquotedblleft{}distance\textquotedblright{} and \emph{\ru{далёкий}} \textquotedblleft{}far\textquotedblright{}: to move a thing far away.
\end{cousinweb}

\begin{grammarlens}[title={What you've built}]
One more word for telling, explaining and describing at B1.
\end{grammarlens}

\subsection*{Your turn}
\begin{itemize}
\raggedright
  \item \emph{Say it:} \emph{udalyát}
  \item \emph{Say it:} \emph{udalyát}, once more
  \item \emph{Say it:} a sentence of your own with \emph{udalyát}
  \item \emph{Recall:} say the Russian for to install, then the Russian for to update, then say \emph{udalyát} again
\end{itemize}

\begin{tcolorbox}[breakable,colback=teal!4,colframe=teal!35!black,title={Before you move on}]
What does \ru{удалять} mean? (\textquotedblleft{}To delete, to remove\textquotedblright{}.) Say it once more.
\end{tcolorbox}
\section[\texorpdfstring{\ru{загружать} (zagruzhát)}{загружать (zagruzhát)}]{\ru{загружать} (zagruzhát) — to load, to upload; to keep busy}

\label{lesson:RU-C401-zagruzhat}



\begin{quote}
Before the new one: say the Russian for to update, then the Russian for to delete.
\end{quote}

\subsection*{You'll want to know: \ru{загружать}}
\textbf{\ru{загружать}} (\emph{zagruzhát}) — \textquotedblleft{}to load, to upload; to keep busy\textquotedblright{}.

\begin{quote}
\ru{Страница} \ru{долго} \ru{загружается}.
\end{quote}

\begin{quote}
\emph{Stranítsa dólgo zagruzháyetsya.}
\end{quote}

\begin{quote}
The page takes a long time to load.
\end{quote}

Imperfective; perfective \emph{\ru{загрузить}}. Of work: \emph{\ru{меня} \ru{загрузили} \ru{работой}} — they have loaded me with work.

\begin{cousinweb}
\textbf{\ru{загружать}} is built from \emph{\ru{за}-} and \emph{\ru{груз}} \textquotedblleft{}a load\textquotedblright{}, the root of \emph{\ru{грузовик}} \textquotedblleft{}a lorry\textquotedblright{}: to put a load in.
\end{cousinweb}

\begin{grammarlens}[title={What you've built}]
One more word for telling, explaining and describing at B1.
\end{grammarlens}

\subsection*{Your turn}
\begin{itemize}
\raggedright
  \item \emph{Say it:} \emph{zagruzhát}
  \item \emph{Say it:} \emph{zagruzhát}, once more
  \item \emph{Say it:} a sentence of your own with \emph{zagruzhát}
  \item \emph{Recall:} say the Russian for to update, then the Russian for to delete, then say \emph{zagruzhát} again
\end{itemize}

\begin{tcolorbox}[breakable,colback=teal!4,colframe=teal!35!black,title={Before you move on}]
What does \ru{загружать} mean? (\textquotedblleft{}To load, to upload; to keep busy\textquotedblright{}.) Say it once more.
\end{tcolorbox}
